% Pursuing Stacks (cotes 134-2 to 134-8): archive cote 134-6, site page 101 (PDF page 102), typescript page 419.
% typescript. Draft compared: ps5.tex:6734-6808.
% AI-transcribed by Claude Opus 5.5 (claude-opus-5-5, adaptive thinking, effort high), Message Batches API, batch ps3 (msgbatch_01XdFmuVVBAcgc7EFxdtCqsj), 2026-11-12; not yet checked by eye.
\scanpage{134-6/101}
\mspage{419}

(90) \qquad $(L',F)\hand{$\longmapsto$}\underline{\mathrm{Hom}}(L',F) \ :\ B^{\wedge}_{ab} \times A^{\wedge}_{\hand{$\mathcal{N}$}} \hand{$\longrightarrow$} A^{\wedge}_{\hand{$\mathcal{N}$}}$

is defined when the additive category \hand{$\mathcal{N}$} is stable under small inverse
limits, and can be equally described as taking external Hom's component-%
wise
\[ \underline{\mathrm{Hom}}(L',F)\ (a) \ = \ \mathrm{Hom}(L'(a),F(a)) \qquad . \]
These operations make a sense too when $A^{\wedge}$ is replaced by an arbitrary topos
\hand{$\mathcal{A}$} , $B^{\wedge}_{ab}$ being replaced by the category of abelian cosheaves on \hand{$\mathcal{A}$}. It
does'nt seem worthwhile here to dwell on them, as they do'nt seem to be
so relevant for the homology and cohomology formalism we want to develop
in the next sections. I like to point out though that in the cohomology
formalism of ringed topoi the tensor product operation (89) and the \xout{dxxx}
derived operation $\overset{L}{\otimes}$ on the relevant derived categories\inserted{$D_{\bullet}$} play an important
role, and \inserted{it} is likely therefore that in a more extensive development of
homology and cohomology formalism within the context of topoi $A^{\wedge}$ and maps
\margin{\unsure{paragraph sign}} in (Cat), the same will hold for the dual operation (90) \inserted{too.}

The reader
who may feel confused by the manifold use of the symbol \uline{Hom} should notice
that there is no possibility of confusion reasonably between (90) and
(65) (p.409), as the argument $L'$ in (90) is in $B^{\wedge}_{ab}$ , whereas the argu-%
ment $L$ in (65) is in an altogether different category $A^{\wedge}_{ab}$ . In the case
when $A$ is the final category say hence $A=B$ and a confusion might arise,
the two operations turn out to be actually the same (up to canonical iso-%
morphism). A similar remark applies to the fear of confusion between the
kindred operations $\otimes$ and $\circledast$. I daresay I devoted a considerable
amount of attention on terminology and notation around the abelianization
story - and it does seem that a pretty coherent formalism is emerging
indeed.

17.8.

107. \quad F) \quad \uline{Extension of ground ring from} $\underline{Z}$ \uline{to} $k$\unsure{$_{*}$} (k-linearization).

I would like still to make a quick review of the main facts and formu-%
las of the last two sections, replacing throughout the ground ring $\underline{Z}$
by an arbitrary \uline{commutative} ring k , \deleted{replacing throughout} \inserted{and} additive catego-%
ries \hand{$\mathcal{M}$} and additive functors between these, by k-additive categories and
k-additive functors. This will allow us to check that the conceptual and
notational set-up we got so far extends smoothly to k-linearization.

\margin{101}
